\clearpage\begingroup\SourceLeaf{111}\HistoricalHead{106}{L}{PHYSIK. BRIEFWECHSEL.}
\begin{center}[20.]\end{center}
\begin{center}\textsc{Gauss} an \textsc{Schumacher}.\quad Göttingen, 17. Januar 1836.\end{center}
\begin{center}\begin{minipage}{.82\linewidth}\centering\hrule\vspace{3pt}
\small [Briefwechsel zwischen \textsc{Gauss} und \textsc{Schumacher} II, S. 435.]
\par\vspace{3pt}\hrule\end{minipage}\end{center}
\noindent .\ .\ .\ .\ .\ .\ .\ .\ .\ .\ .\ .\ .\ .\ .\ .\ .\ .\ .\ .\ .\par
\indent Aus dem reichen Gebiet des Electromagnetismus habe ich seit Anfang
dieses Jahrs noch eine schöne Provinz unsern Apparaten unterworfen, nemlich
die Thermogalvanische Erregung. Vermittelst einer besondern einfachen Vorrichtung
bringe ich auch \textls[120]{diese} Ströme in solcher Stärke hervor, dass von
der Sternwarte aus eine 24pfündige Nadel im physikalischen Cabinett in grosse
Bewegung gesetzt wird. Ich denke, dies wird ein höchst wichtiges Mittel sein,
um in der physikalischen Theorie der Wärme die interessantesten und feinsten
Experimente zu machen. — Einige andere damit verwandte Versuche, habe
ich mir auf die nächste Zeit vorgesetzt. \mbox{.\ .\ .\ .\ .\ .\ .\ .\ .\ .\ .\ .\ .\ .\ .\ .\ .\ .}

\bigskip\Needspace{14\baselineskip}
\begin{center}[21.]\end{center}
\begin{center}\textsc{Gauss} an \textsc{Olbers}.\quad Göttingen, 1. März 1836.\end{center}
\begin{center}\begin{minipage}{.88\linewidth}\centering\hrule\vspace{3pt}
\small [\textsc{Wilhelm Olbers}, Sein Leben und seine Werke II, 2, S. 633.]
\par\vspace{3pt}\hrule\end{minipage}\end{center}
\noindent .\ .\ .\ .\ .\ .\ .\ .\ .\ .\ .\ .\ .\ .\ .\ .\ .\ .\ .\ .\ .\par
\indent Seit Jan[uar] habe ich auch einige Versuche über den sogenannten Thermomagnetismus,
richtiger Thermogalvanismus angefangen; wir haben neue
Apparate angeordnet, womit man die Wirkung ganz ausserordentlich verstärken
kann. Vielleicht kann man davon für die Pyrometrie (Hochöfen,
Porzellanöfen, Schmelzöfen etc.) sehr wichtige Anwendungen machen. Meine
eigenen Versuche bleiben freilich bei den mesquinen, mir zu Gebote stehenden
Mitteln immer nur auf einen kleinen Maassstab beschränkt. \mbox{.\ .\ .\ .\ .\ .\ .\ .\ .\ .}
\clearpage\endgroup
